\documentclass[11pt,a4paper]{article}

\usepackage[margin=2.3cm]{geometry}
\usepackage{amsmath,amssymb,bm}
\usepackage{booktabs}
\usepackage{array}
\usepackage{longtable}
\usepackage{enumitem}
\usepackage{xcolor}
\usepackage{hyperref}
\usepackage{listings}
\usepackage{microtype}

\hypersetup{
    colorlinks=true,
    linkcolor=blue,
    urlcolor=blue,
    citecolor=blue
}

\lstset{
    basicstyle=\ttfamily\small,
    breaklines=true,
    frame=single,
    columns=fullflexible,
    keepspaces=true
}

\setlength{\parindent}{2em}
\setlength{\parskip}{0.4em}
\renewcommand{\arraystretch}{1.25}

\title{Notes on On-Policy Distillation}
\author{}
\date{}

\begin{document}

\maketitle
\tableofcontents
\newpage

\section{Notation}

\begin{itemize}[leftmargin=2em]
    \item $x=(x_1,\ldots,x_T)$: the complete sequence generated by the student.
    \item $x_{1..t}$: the prefix after generating the first $t$ tokens.
    \item $x_{t+1}$: the next token.
    \item $\pi_\theta$: the student model.
    \item $\pi_{\text{teacher}}$: the fixed teacher model.
    \item $\theta$: the student-model parameters.
    \item $\gamma$: the discount factor for future KL.
    \item $\beta$: the code's \texttt{kl\_penalty\_coef}.
\end{itemize}

Two additional symbols are used for importance sampling:

\begin{itemize}[leftmargin=2em]
    \item $q$: the sampler policy that actually generates the rollout.
    \item $p_\theta$: the learner policy used during the training forward pass.
\end{itemize}

In the ideal strict on-policy setting, both should be identical to $\pi_\theta$. In a real system, however, there is usually a mismatch between rollout and training: for example, the training and inference infrastructure may differ, and when the same rollout is reused for multiple parameter updates, the learner gradually drifts away from the sampler that generated the data. In addition, some infrastructure now uses importance sampling to transmit the RL gradient.

\section{Starting from the Per-Token Reverse KL}

For a fixed prefix $x_{1..t}$, define the reverse KL from the student to the teacher over the next token:

\[
\boxed{
\mathrm{KL}\left(\pi_\theta\|\pi_{\text{teacher}}\right)
=
\mathbb E_{x\sim\pi_\theta}
\left[
\log \pi_\theta(x_{t+1}|x_{1..t})
-
\log \pi_{\text{teacher}}(x_{t+1}|x_{1..t})
\right]
}
\]

Expanding it more precisely as a sum over the full vocabulary gives:

\[
\begin{aligned}
D_t(x_{1..t})
&=
\sum_{x_{t+1}}
\pi_\theta(x_{t+1}|x_{1..t})\\
&\qquad\cdot
\left[
\log\pi_\theta(x_{t+1}|x_{1..t})
-
\log\pi_{\text{teacher}}(x_{t+1}|x_{1..t})
\right].
\end{aligned}
\]

It compares the two next-token conditional distributions under the same prefix visited by the student:

\[
\pi_\theta(\cdot|x_{1..t})
\quad\text{and}\quad
\pi_{\text{teacher}}(\cdot|x_{1..t}).
\]

For simpler notation, define the sampled single-token KL:

\[
d_t
=
\log\pi_\theta(x_{t+1}|x_{1..t})
-
\log\pi_{\text{teacher}}(x_{t+1}|x_{1..t}).
\]

\section{How the Full-Vocabulary Reverse KL Produces the Final Gradient}

Fix $x_{1..t}$ and differentiate $D_t$:

\[
\begin{aligned}
\nabla_\theta D_t
&=
\sum_{x_{t+1}}
\nabla_\theta\pi_\theta(x_{t+1}|x_{1..t})\,d_t\\
&\quad+
\sum_{x_{t+1}}
\pi_\theta(x_{t+1}|x_{1..t})
\nabla_\theta\log\pi_\theta(x_{t+1}|x_{1..t}).
\end{aligned}
\]

Using the log-derivative identity

\[
\nabla_\theta\pi_\theta
=
\pi_\theta\nabla_\theta\log\pi_\theta,
\]

we obtain:

\[
\nabla_\theta D_t
=
\mathbb E_{x_{t+1}\sim\pi_\theta}
\left[
(d_t+1)
\nabla_\theta\log\pi_\theta(x_{t+1}|x_{1..t})
\right].
\]

Moreover:

\[
\begin{aligned}
\mathbb E_{\pi_\theta}
\left[\nabla_\theta\log\pi_\theta\right]
&=
\sum_{x_{t+1}}
\pi_\theta(x_{t+1}|x_{1..t})
\nabla_\theta\log\pi_\theta(x_{t+1}|x_{1..t})\\
&=
\sum_{x_{t+1}}
\nabla_\theta\pi_\theta(x_{t+1}|x_{1..t})\\
&=\nabla_\theta 1\\
&=0.
\end{aligned}
\]

Therefore, the $+1$ disappears in expectation:

\[
\boxed{
\nabla_\theta D_t
=
\mathbb E_{x_{t+1}\sim\pi_\theta}
\left[
d_t
\nabla_\theta\log\pi_\theta(x_{t+1}|x_{1..t})
\right]
}
\]

Equivalently:

\[
\boxed{
\nabla_\theta D_t
=
\mathbb E_{x_{t+1}\sim\pi_\theta}
\left[
\left(
\log\pi_\theta(x_{t+1}|x_{1..t})
-
\log\pi_{\text{teacher}}(x_{t+1}|x_{1..t})
\right)
\nabla_\theta\log\pi_\theta(x_{t+1}|x_{1..t})
\right].
}
\]

This is the core token-level policy-gradient form used by OPD.

\section{Estimating the Full-Vocabulary Gradient with Sampled Tokens}

Because computing the reverse KL over the full vocabulary is expensive, and currently only DeepSeek v4 is known to use it, practitioners generally use sampled tokens or top-$K$ approximations. We first show that the Monte Carlo average over $K$ sampled tokens is unbiased, and then consider the commonly used special case $K=1$.

The previous section gives:

\[
\boxed{
\nabla_\theta D_t
=
\mathbb E_{x_{t+1}\sim\pi_\theta}
\left[
d_t\nabla_\theta\log\pi_\theta(x_{t+1}|x_{1..t})
\right].
}
\]

Expanding this expectation as a full-vocabulary sum:

\[
\begin{aligned}
\nabla_\theta D_t
&=
\sum_{x_{t+1}}
\pi_\theta(x_{t+1}|x_{1..t})\\
&\qquad\cdot
\left[
\left(
\log\pi_\theta(x_{t+1}|x_{1..t})
-
\log\pi_{\text{teacher}}(x_{t+1}|x_{1..t})
\right)
\nabla_\theta\log\pi_\theta(x_{t+1}|x_{1..t})
\right].
\end{aligned}
\]

Thus, the true gradient enumerates every possible next token and uses the student probability

\[
\pi_\theta(x_{t+1}|x_{1..t})
\]

as the weight in the sum.

Instead of enumerating the full vocabulary, sample $K$ tokens from the student distribution:

\[
x_{t+1}^{(1)},\ldots,x_{t+1}^{(K)}
\sim
\pi_\theta(\cdot|x_{1..t}),
\]

where sample $k$ has:

\[
d_t^{(k)}
=
\log\pi_\theta(x_{t+1}^{(k)}|x_{1..t})
-
\log\pi_{\text{teacher}}(x_{t+1}^{(k)}|x_{1..t}).
\]

For sample $k$, construct the single-sample gradient estimator:

\[
\hat g^{(k)}
=
d_t^{(k)}
\nabla_\theta
\log\pi_\theta(x_{t+1}^{(k)}|x_{1..t}).
\]

Average over the $K$ samples:

\[
\boxed{
\hat g_K
=
\frac1K
\sum_{k=1}^{K}
\hat g^{(k)}
=
\frac1K
\sum_{k=1}^{K}
d_t^{(k)}
\nabla_\theta
\log\pi_\theta(x_{t+1}^{(k)}|x_{1..t}).
}
\]

For a single sample:

\[
\begin{aligned}
\mathbb E[\hat g^{(k)}]
&=
\sum_{x_{t+1}^{(k)}}
\pi_\theta(x_{t+1}^{(k)}|x_{1..t})\\
&\qquad\cdot
\left[
\left(
\log\pi_\theta(x_{t+1}^{(k)}|x_{1..t})
-
\log\pi_{\text{teacher}}(x_{t+1}^{(k)}|x_{1..t})
\right)
\nabla_\theta
\log\pi_\theta(x_{t+1}^{(k)}|x_{1..t})
\right]\\
&=\nabla_\theta D_t.
\end{aligned}
\]

Therefore:

\[
\begin{aligned}
\mathbb E[\hat g_K]
&=
\frac1K\sum_{k=1}^{K}
\mathbb E
\left[
d_t^{(k)}
\nabla_\theta\log\pi_\theta(x_{t+1}^{(k)}|x_{1..t})
\right]\\
&=\frac1K\sum_{k=1}^{K}\nabla_\theta D_t\\
&=\nabla_\theta D_t.
\end{aligned}
\]

Thus, the average over $K$ sampled tokens is an unbiased Monte Carlo estimator of the full-vocabulary gradient.

In practice, $K=1$ is common. The estimator becomes:

\[
\boxed{
\hat g_1
=
d_t\nabla_\theta\log\pi_\theta(x_{t+1}|x_{1..t}),
\qquad
x_{t+1}\sim\pi_\theta.
}
\]

It remains unbiased, but its variance is higher than that of an estimator using a larger $K$.

There is an important engineering detail missing from the derivation above: a sampled-token estimator cannot simply treat the sampled KL as an ordinary loss and backpropagate through it; it requires stop-gradient.

I initially overlooked this as well, until the official code made the issue apparent. After sampling,

\[
x_{t+1}\sim\pi_\theta,
\]

automatic differentiation treats $x_{t+1}$ as a fixed token; the sampling probability itself does not remain explicitly in the computation graph. The implementation therefore needs to reproduce the score-function gradient estimator from the previous section:

\[
\hat g_K
=
\frac1K
\sum_{k=1}^{K}
d_t^{(k)}
\nabla_\theta
\log\pi_\theta(x_{t+1}^{(k)}|x_{1..t}).
\]

The surrogate loss is:

\[
\boxed{
\mathcal L_{\mathrm{MC}}
=
\frac1K
\sum_{k=1}^{K}
\operatorname{sg}\!\left[d_t^{(k)}\right]
\log\pi_\theta(x_{t+1}^{(k)}|x_{1..t})
}
\]

where $\operatorname{sg}$ denotes stop-gradient. It treats $d_t^{(k)}$ only as a numerical weight and lets automatic differentiation pass through the score function:

\[
\nabla_\theta\mathcal L_{\mathrm{MC}}
=
\frac1K
\sum_{k=1}^{K}
\operatorname{sg}\!\left[d_t^{(k)}\right]
\nabla_\theta
\log\pi_\theta(x_{t+1}^{(k)}|x_{1..t}).
\]

Without stop-gradient, the expression $d_t\log\pi_\theta(x_{t+1}|x_{1..t})$ would produce an additional term:

\[
\log\pi_\theta(x_{t+1}|x_{1..t})\nabla_\theta d_t.
\]

The difference from the full-vocabulary KL is that the full objective

\[
\mathcal L_t^{\text{full}}
=
\sum_{x_{t+1}}
\pi_\theta(x_{t+1}|x_{1..t})
\left[
\log\pi_\theta(x_{t+1}|x_{1..t})
-
\log\pi_{\text{teacher}}(x_{t+1}|x_{1..t})
\right]
\]

explicitly includes the probability weight $\pi_\theta(x_{t+1}|x_{1..t})$ for every token in the sum. It therefore differentiates the complete true objective directly and does not require this special stop-gradient treatment.

\section{Sampled-Token Loss versus Full-Vocabulary Loss}

Both methods estimate the same conditional KL for a fixed prefix, but they use different implementations.

\begin{longtable}{>{\raggedright\arraybackslash}p{0.22\textwidth}
                  >{\raggedright\arraybackslash}p{0.34\textwidth}
                  >{\raggedright\arraybackslash}p{0.34\textwidth}}
\toprule
Aspect & Full-vocabulary loss & Sampled-token loss \\
\midrule
Next token & Enumerate every vocabulary token & Use only the actually sampled token \\
Objective & Compute the conditional KL exactly & Monte Carlo estimate of the conditional-KL gradient \\
Teacher output & Full-vocabulary probabilities or logits & Logprob of the sampled token \\
Variance & No sampling variance for the current prefix & Relatively high when $K=1$ \\
Compute & Scales with vocabulary size & Only the teacher logprob of sampled tokens is needed per position \\
Autodiff & Backpropagate through the full KL directly & Use a policy-gradient surrogate \\
Stop-gradient & No special stop-gradient required & Required when $d_t$ is used as an advantage \\
\bottomrule
\end{longtable}

\section{Future-Token Effects and the Discount Factor}

So far, we have analyzed the next-token conditional KL under a fixed prefix $x_{1..t}$. A common question is: when the current token changes, the prefixes visited later also change; should the KL at those future positions be credited to the current token?

If the entire rollout is viewed as a sample generated by the student, a natural sequence-level objective is the sum of the KL terms along the trajectory:

\[
\mathcal J(\theta)
=
\mathbb E_{x\sim\pi_\theta}
\left[\sum_{t=0}^{T-1}d_t\right],
\]

where:

\[
d_t
=
\log\pi_\theta(x_{t+1}|x_{1..t})
-
\log\pi_{\text{teacher}}(x_{t+1}|x_{1..t}).
\]

This can also be understood as the reverse KL between the complete sequence distributions, because:

\[
\log\pi_\theta(x)
=
\sum_t\log\pi_\theta(x_{t+1}|x_{1..t}),
\]

and likewise for the teacher.

Differentiating $\mathcal J$ gives:

\[
\begin{aligned}
\nabla_\theta\mathcal J
=
\mathbb E_{x\sim\pi_\theta}
\Bigg[
&\left(\sum_t\nabla_\theta
\log\pi_\theta(x_{t+1}|x_{1..t})\right)
\left(\sum_k d_k\right)\\
&+\sum_k
\nabla_\theta\log\pi_\theta(x_{k+1}|x_{1..k})
\Bigg].
\end{aligned}
\]

The expectation of the final term is zero. Using causality, the score at time $t$ does not need to multiply past $d_k$ terms with $k<t$, yielding:

\[
\boxed{
\nabla_\theta\mathcal J
=
\mathbb E_{x\sim\pi_\theta}
\left[
\sum_{t=0}^{T-1}
\left(\sum_{k=t}^{T-1}d_k\right)
\nabla_\theta\log\pi_\theta(x_{t+1}|x_{1..t})
\right].
}
\]

Therefore, in the strict sequence-level reverse-KL gradient, the weight for token $t$ should be:

\[
d_t+d_{t+1}+\cdots+d_{T-1}.
\]

The reason is that the current token affects not only the current KL, but also every later prefix visited by the rollout.

We therefore define the discounted future KL:

\[
G_t^{(\gamma)}
=
\sum_{k=t}^{T-1}\gamma^{k-t}d_k.
\]

That is:

\[
G_t^{(\gamma)}
=
d_t+\gamma d_{t+1}+\gamma^2d_{t+2}+\cdots.
\]

The corresponding KL advantage is:

\[
A_t^{(\gamma)}=-\beta G_t^{(\gamma)}.
\]

The meaning of different $\gamma$ values is as follows.

\subsection{\texorpdfstring{$\gamma=1$}{gamma=1}: Full Sequence Credit Assignment}

\[
G_t^{(1)}=d_t+d_{t+1}+\cdots+d_{T-1}.
\]

This corresponds to undiscounted credit assignment under the sequence-level objective.

\subsection{\texorpdfstring{$0<\gamma<1$}{0<gamma<1}: Partial Future Credit Assignment}

\[
G_t^{(\gamma)}
=d_t+\gamma d_{t+1}+\gamma^2d_{t+2}+\cdots.
\]

Future effects are retained, but more distant KL terms receive smaller weights.

\subsection{\texorpdfstring{$\gamma=0$}{gamma=0}: Immediate Token Only}

\[
\boxed{G_t^{(0)}=d_t.}
\]

The current token is updated only from the teacher signal at its own position and does not take responsibility for future KL. This is also the current default. The main reason is that the TML official blog mentions that a positive discount factor is mathematically more complete, but experiments did not show a performance improvement; for simplicity, zero discount is used deliberately.

\section{The Final Unified OPD Formula}

\subsection{Rollout}

\[
x_{t+1}\sim q(\cdot|x_{1..t}).
\]

\subsection{Sampled Reverse KL}

\[
d_t
=
\log q(x_{t+1}|x_{1..t})
-
\log\pi_{\text{teacher}}(x_{t+1}|x_{1..t}).
\]

\subsection{Discounted Future KL}

\[
G_t^{(\gamma)}
=
\sum_{k=t}^{T-1}\gamma^{k-t}d_k.
\]

\subsection{Advantage}

\[
\boxed{
A_t
=
-\beta\operatorname{stopgrad}
\left[G_t^{(\gamma)}\right].
}
\]

\subsection{Importance Ratio}

\[
\boxed{
\rho_t(\theta)
=
\frac{
p_\theta(x_{t+1}|x_{1..t})
}{
q(x_{t+1}|x_{1..t})
}.
}
\]

\subsection{Tinker Loss}

\[
\boxed{
\mathcal L_{\mathrm{OPD}}
=
-\sum_t\rho_t(\theta)A_t.
}
\]

\subsection{Gradient}

\[
\boxed{
\nabla_\theta\mathcal L_{\mathrm{OPD}}
=
\beta\sum_t
\rho_t(\theta)G_t^{(\gamma)}
\nabla_\theta\log p_\theta(x_{t+1}|x_{1..t}).
}
\]

By default:

\[
\gamma=0,
\]

so:

\[
G_t^{(0)}=d_t,
\]

and the final expression is:

\[
\boxed{
\nabla_\theta\mathcal L_{\mathrm{OPD}}
=
\beta\sum_t
\frac{p_\theta(x_{t+1}|x_{1..t})}
{q(x_{t+1}|x_{1..t})}
\left[
\log q(x_{t+1}|x_{1..t})
-
\log\pi_{\text{teacher}}(x_{t+1}|x_{1..t})
\right]
\nabla_\theta\log p_\theta(x_{t+1}|x_{1..t}).
}
\]

Under the strict on-policy setting,

\[
p_\theta=q=\pi_\theta,
\]

which simplifies to:

\[
\boxed{
\nabla_\theta\mathcal L_{\mathrm{OPD}}
=
\beta\sum_t
\left[
\log\pi_\theta(x_{t+1}|x_{1..t})
-
\log\pi_{\text{teacher}}(x_{t+1}|x_{1..t})
\right]
\nabla_\theta\log\pi_\theta(x_{t+1}|x_{1..t}).
}
\]

\end{document}
